
\documentclass{article}
\usepackage{colortbl}
\usepackage{makecell}
\usepackage{multirow}
\usepackage{supertabular}

\begin{document}

\newcounter{utterance}

\centering \large Interaction Transcript for game `mmlu\_pro', experiment `mmlu\_pro\_math', episode 73 with reflexion{-}Qwen3.5{-}9B.
\vspace{24pt}

{ \footnotesize  \setcounter{utterance}{1}
\setlength{\tabcolsep}{0pt}
\begin{supertabular}{c@{$\;$}|p{.15\linewidth}@{}p{.15\linewidth}p{.15\linewidth}p{.15\linewidth}p{.15\linewidth}p{.15\linewidth}}
    \# & \multicolumn{2}{c}{Player} && \multicolumn{2}{c}{Game Master} \\
    \hline

    \theutterance \stepcounter{utterance}  
    & & & \multicolumn{4}{p{0.6\linewidth}}{
        \cellcolor[rgb]{0.9,0.9,0.9}{
            \makecell[{{p{\linewidth}}}]{
                \texttt{\tiny{[P1$\langle$GM]}}
                \texttt{You have attempted this same task before and did not succeed. These are the plans you wrote after those attempts.} \\
\texttt{Trial \#1: New plan: I failed because I guessed (C) for the final question without performing the Newton{-}Raphson calculation, likely relying on an incorrect intuition or random selection. The function \$f(x) = 2x\^{}3 + e\^{}x {-} 10\$ is strictly increasing, and \$f(1) \textbackslash{}approx {-}5.28\$ while \$f(2) \textbackslash{}approx 10.38\$, placing the root between 1 and 2. Option (C) 2.89 is too high (\$f(2.89) \textbackslash{}gg 10\$), whereas Option (F) 1.42 is the only value in the valid interval that approximates the root (\$f(1.42) \textbackslash{}approx 0\$). Next time, I will verify the mathematical validity of each option against the problem constraints before selecting an answer, ensuring the chosen value actually satisfies the equation within a reasonable margin of error.} \\
\\ 
\\ 
\texttt{The following are multiple choice questions (with answers) about math. Answer only with "The answer is (X)" where X is the correct letter choice.} \\
\\ 
\texttt{Question:} \\
\texttt{The symmetric group \$S\_n\$ has \$} \\
\texttt{\textbackslash{}factorial\{n\}\$ elements, hence it is not true that \$S\_\{10\}\$ has 10 elements.} \\
\texttt{Find the characteristic of the ring 2Z.} \\
\texttt{Options:} \\
\texttt{A. 0} \\
\texttt{B. 30} \\
\texttt{C. 3} \\
\texttt{D. 10} \\
\texttt{E. 12} \\
\texttt{F. 50} \\
\texttt{G. 2} \\
\texttt{H. 100} \\
\texttt{I. 20} \\
\texttt{J. 5} \\
\texttt{The answer is (A)} \\
\\ 
\texttt{Question:} \\
\texttt{Let V be the set of all real polynomials p(x). Let transformations T, S be defined on V by T:p(x) {-}> xp(x) and S:p(x) {-}> p'(x) = d/dx p(x), and interpret (ST)(p(x)) as S(T(p(x))). Which of the following is true?} \\
\texttt{Options:} \\
\texttt{A. ST + TS is the identity map of V onto itself.} \\
\texttt{B. TS = 0} \\
\texttt{C. ST = 1} \\
\texttt{D. ST {-} TS = 0} \\
\texttt{E. ST = T} \\
\texttt{F. ST = 0} \\
\texttt{G. ST = TS} \\
\texttt{H. ST {-} TS is the identity map of V onto itself.} \\
\texttt{I. TS = T} \\
\texttt{J. ST = S} \\
\texttt{The answer is (H)} \\
\\ 
\texttt{Question:} \\
\texttt{Let A be the set of all ordered pairs of integers (m, n) such that 7m + 12n = 22. What is the greatest negative number in the set B = \{m + n : (m, n) \textbackslash{}in A\}?} \\
\texttt{Options:} \\
\texttt{A. {-}5} \\
\texttt{B. 0} \\
\texttt{C. {-}3} \\
\texttt{D. {-}7} \\
\texttt{E. {-}4} \\
\texttt{F. {-}6} \\
\texttt{G. {-}1} \\
\texttt{H. {-}2} \\
\texttt{I. {-}9} \\
\texttt{J. N/A} \\
\texttt{The answer is (E)} \\
\\ 
\texttt{Question:} \\
\texttt{A tank initially contains a salt solution of 3 grams of salt dissolved in 100 liters of water. A salt solution containing 0.02 grams of salt per liter of water is sprayed into the tank at a rate of 4 liters per minute. The sprayed solution is continually mixed with the salt solution in the tank, and the mixture flows out of the tank at a rate of 4 liters per minute. If the mixing is instantaneous, how many grams of salt are in the tank after 100 minutes have elapsed?} \\
\texttt{Options:} \\
\texttt{A. 3 + e\^{}{-}2} \\
\texttt{B. 2 {-} e\^{}{-}4} \\
\texttt{C. 2 {-} e\^{}{-}2} \\
\texttt{D. 3 + e\^{}{-}4} \\
\texttt{E. 2 + e\^{}{-}3} \\
\texttt{F. 2 {-} e\^{}{-}3} \\
\texttt{G. 3 {-} e\^{}{-}2} \\
\texttt{H. 2 + e\^{}{-}2} \\
\texttt{I. 2 + e\^{}{-}4} \\
\texttt{J. 2} \\
\texttt{The answer is (I)} \\
\\ 
\texttt{Question:} \\
\texttt{A total of 30 players will play basketball at a park. There will be exactly 5 players on each team. Which statement correctly explains how to find the number of teams needed?} \\
\texttt{Options:} \\
\texttt{A. Multiply 5 by 5 to find 25 teams.} \\
\texttt{B. Divide 30 by 5 to find 6 teams.} \\
\texttt{C. Add 5 to 30 to find 35 teams.} \\
\texttt{D. Subtract 30 from 5 to find {-}25 teams.} \\
\texttt{E. Divide 5 by 30 to find 0.1667 teams.} \\
\texttt{F. Add 5 to 30 then divide by 2 to find 17.5 teams.} \\
\texttt{G. N/A} \\
\texttt{H. N/A} \\
\texttt{I. N/A} \\
\texttt{J. N/A} \\
\texttt{The answer is (B)} \\
\\ 
\texttt{Question:} \\
\texttt{Please solve the equation 2*x\^{}3 + e\^{}x = 10 using newton{-}raphson method.} \\
\texttt{Options:} \\
\texttt{A. 3.14} \\
\texttt{B. 2.05} \\
\texttt{C. 2.89} \\
\texttt{D. {-}1.23} \\
\texttt{E. 2.56} \\
\texttt{F. 1.42} \\
\texttt{G. 0.98} \\
\texttt{H. 0.75} \\
\texttt{I. 1.85} \\
\texttt{J. 3.67} \\
            }
        }
    }
     \\ \\

    \theutterance \stepcounter{utterance}  
    & \multicolumn{4}{p{0.6\linewidth}}{
        \cellcolor[rgb]{0.561,0.737,0.561}{
            \makecell[{{p{\linewidth}}}]{
                \texttt{\tiny{[P1$\rangle$GM]}}
                \texttt{The answer is (F)} \\
            }
        }
    }
    & & \\ \\

    \theutterance \stepcounter{utterance}  
    & & & \multicolumn{2}{p{0.3\linewidth}}{
        \cellcolor[rgb]{0.9,0.9,0.9}{
            \makecell[{{p{\linewidth}}}]{
                \texttt{\tiny{[GM$|$GM]}}
                \texttt{F} \\
            }
        }
    }
    & & \\ \\

    \theutterance \stepcounter{utterance}  
    & & & \multicolumn{2}{p{0.3\linewidth}}{
        \cellcolor[rgb]{0.9,0.9,0.9}{
            \makecell[{{p{\linewidth}}}]{
                \texttt{\tiny{[GM$|$GM]}}
                \texttt{F} \\
            }
        }
    }
    & & \\ \\

    \theutterance \stepcounter{utterance}  
    & & & \multicolumn{2}{p{0.3\linewidth}}{
        \cellcolor[rgb]{0.9,0.9,0.9}{
            \makecell[{{p{\linewidth}}}]{
                \texttt{\tiny{[GM$|$GM]}}
                \texttt{game\_result = WIN} \\
            }
        }
    }
    & & \\ \\

\end{supertabular}
}

\end{document}
